% Part: proof-theory
% Chapter: sequent-calculus
% Section: admissible-derivable

\documentclass[../../../include/open-logic-section]{subfiles}

\begin{document}

\olfileid{pt}{seq}{adm}

\olsection{可容许规则与可导出规则}

证明论的许多结果涉及——或者说使用了——关于是否所有使用特定规则可导出的内容，在没有该规则时也可导出的结论。切割消去定理就是最著名的例子。它表明，任何能够在一个带有~\Cut{} 规则的相继式系统中推导出的相继式，也能够在没有~\Cut{} 的情况下推导出来。由于这类性质非常重要，它有自己的名称。

\begin{defn}
在一个相继式系统中，如果任何在该系统连同规则~$R$ 中能够被证明的相继式，不使用~$R$ 也能被证明，则称规则~$R$ 是\emph{可容许的}。
\end{defn}

\begin{ex}\ollabel{ex:land-deriv}
$\LeftR{\land}$ 规则在 \Log{G1c} 和 \Log{G3c} 中具有不同的形式。在 \Log{G1c} 中，该规则有两个版本：一个在前提中包含主公式~$!A \land !B$ 的左合取支 $!A$，另一个在前提中包含右合取支~$!B$：
\[
\Axiom$ !A, \Gamma \fCenter \Delta$
\RightLabel{$\LeftR{\land}_1$}
\UnaryInf$ !A \land !B, \Gamma \fCenter \Delta$
\DisplayProof
\qquad
\Axiom$!B, \Gamma \fCenter \Delta$
\RightLabel{$\LeftR{\land}_2$}
\UnaryInf$!A \land !B, \Gamma \fCenter \Delta$
\DisplayProof
\]
相比之下，\Log{G3c} 只有一个规则：
\[
\Axiom$ !A, !B, \Gamma \fCenter \Delta$
\RightLabel{$\LeftR{\land}_3$}
\UnaryInf$ !A \land !B, \Gamma \fCenter \Delta$
\DisplayProof
\]
我们现在可能会问：“\Log{G3c} 的规则 $\LeftR{\land}_3$ 在 \Log{G1c} 中是可容许的吗？”答案是肯定的：我们可以仅使用~\Log{G1c} 的规则来直接模拟任何使用 $\LeftR{\land}_3$ 的推理。除了 $\LeftR{\land}_1$ 和 $\LeftR{\land}_2$，我们还需要结构规则~$\LeftR{\Contraction}$：
\[
\Axiom$!A, !B, \Gamma \fCenter \Delta$
\RightLabel{$\LeftR{\land}_1$}
\UnaryInf$!A \land !B, !B, \Gamma \fCenter \Delta$
\RightLabel{$\LeftR{\land}_2$}
\UnaryInf$!A \land !B, !A \land !B, \Gamma \fCenter \Delta$
\RightLabel{\LeftR{\Contraction}}
\UnaryInf$!A \land !B \Gamma \fCenter \Delta$
\DisplayProof
\]
\end{ex}

使用其他规则来模拟某条规则所作的推理，是证明该规则可容许的一种方法。但这并不是唯一的方法，有些规则是可容许的，但却不能如此被模拟。能够如此模拟的规则被称为\emph{可导出的}。

\begin{defn}
一个规则~$R$，
\[
\AxiomC{$S_1$}
\AxiomC{$\dots$}
\AxiomC{$S_n$}
\RightLabel{$R$}
\TrinaryInfC{$S$}
\DisplayProof
\]
在一个系统中被称为\emph{可导出的}，如果存在相继式~$S$ 的一个图式推导，该推导使用该系统的公理和规则，并将形如~$S_1$, \dots,~$S_n$ 的相继式作为额外初始相继式。
\end{defn}

在 \olref{ex:land-deriv} 中给出的推导表明，$\LeftR{\land}_3$ 是~\Log{G1c} 中的一个可导出规则，因为该推导正是一个图式推导：它将 $\LeftR{\land}_3$ 的前提作为初始相继式，仅使用 \Log{G1c} 的规则推导出 $\LeftR{\land}_3$ 的结论。（它没有使用 \Log{G1c} 的任何公理。）

\begin{prop}\ollabel{prop:lor-G3-adm}
\Log{G3c} 的规则 $\LeftR{\land}$ 和 $\RightR{\lor}$ 在~\Log{G1c} 中是可导出的。
\end{prop}

\begin{proof}
  $\LeftR{\land}$ 的证明已在 \olref{ex:land-deriv} 中给出。$\RightR{\lor}$ 的证明留作练习。
\end{proof}

\begin{prob}
证明 \Log{G3c} 的规则 $\RightR{\lor}$ 在~\Log{G1c} 中是可导出的。
\end{prob}

\begin{prob}
假设 $\liff$ 是一个被定义的逻辑算子：$!A \liff !B$ 是
$(!A \lif !B) \land (!B \lif !A)$
的简写。证明以下规则
\begin{enumerate}
\item \Axiom$\Gamma, !A, !B \fCenter \Delta$
\Axiom$\Gamma \fCenter \Delta, !A, !B$
\RightLabel{\LeftR{\liff}}
\BinaryInf$\Gamma \fCenter \Delta, !A \liff !B$
\DisplayProof
\item
\Axiom$!A, \Gamma \fCenter \Delta, !B$
\Axiom$!B, \Gamma \fCenter \Delta, !A$
\RightLabel{\RightR{\liff}}
\BinaryInf$\Gamma \fCenter \Delta, !A \liff !B$
\DisplayProof
\end{enumerate}
在 (a)~\Log{G1c} 和 (b)~\Log{G3c} 中是可导出的。
\end{prob}

\begin{ex}\ollabel{ex:weak-adm}
\Log{G1c} 的规则 $\RightR{\Weakening}$，
\[
\Axiom$ \Gamma \fCenter \Delta$
\RightLabel{\RightR{\Weakening}}
\UnaryInf$ \Gamma \fCenter \Delta, !A$
\DisplayProof
\]
在~\Log{G3c} 中是可容许的。证明的思路很简单：假设在~\Log{G3c} 中存在 $\Gamma \Sequent \Delta$ 的一个推导~$\pi$。将公式~$!A$ 添加到~$\pi$ 中每个相继式的后件。公理仍然是公理；正确的推理仍保持正确。新推导的结束式就是 $\Gamma \Sequent \Delta, !A$。

然而，规则 $\RightR{\Weakening}$ 在~\Log{G3c} 中不是可导出的。因为假设它是可导出的，即假设能够找到一个 $\Gamma \Sequent \Delta, !A$ 的推导，该推导仅使用~\Log{G3c} 的公理和规则，并可能以~$\Gamma \Sequent \Delta$ 作为额外初始相继式。如果这个推导是纯图式的（这是 $\RightR{\Weakening}$ 成为可推导规则的要求），那么我们可以通过删除 $\Gamma$ 和~$\Delta$ 将它变成一个~$\Sequent !A$ 的推导。这将把额外的初始相继式~$\Gamma \Sequent \Delta$ 变成空相继式~$\quad \Sequent \quad$。由于空相继式不包含任何公式，而~\Log{G3c} 的所有规则都要求其前提中至少有一个活跃公式，因此空相继式不能成为这个图式推导中任何推理的前提。所以，只有当这个推导实际上不使用额外相继式 $\Gamma \Sequent \Delta$ 作为初始相继式时，它才可能是纯图式的。换言之，它将是仅使用~\Log{G3c} 的公理和规则对 $\Gamma \Sequent \Delta, !A$ 所作的图式推导，从而我们就证明了任何这种形式的相继式在~\Log{G3c} 中都有一个推导。但这是不可能的，因为 \Log{G3c} 是可靠的，因此只能推导出在所有结构中都为真的相继式。取 $\Gamma = \Delta = \emptyset$ 且 $!A \ident \lfalse$，\Log{G3c} 将不得不推导出例如相继式 $\quad \Sequent \lfalse$，而它并不能推导出该式。
\end{ex}

下面我们为“弱化在~\Log{G3c} 中是可容许规则”这一结果给出一个漂亮的归纳证明。事实上，正如直观论证所示，一个稍强的结果成立：通过将 $!A$ 加到~$\pi$ 的每个相继式的右边，我们并没有改变推导的结构，特别地，没有增加其高度。由于这一点很重要，它值得拥有自己的定义。

\begin{defn}
一个规则~$R$，
\[
\AxiomC{$S_1$}
\AxiomC{$\dots$}
\AxiomC{$S_n$}
\RightLabel{$R$}
\TrinaryInfC{$S$}
\DisplayProof
\]
在一个系统中是\emph{保持高度可容许的}（或称 \emph{hp-可容许的}），如果只要对 $i = 1$, \dots,~$n$ 有 $\Proves[h_i] S_i$，就有 $\Proves_m S$ 且 $m = \max(h_1, \dots,
h_n)$。
\end{defn}

注意，要使一个规则保持高度可容许，只需：若前提~$S_i$ 有高度为~$h_i = \pheight{\pi_i}$ 的推导~$\pi_i$，则存在结论的一个推导~$\pi$，其高度~$\pheight{\pi}$ 不超过~$h_i$ 的最大值即可。特别地，如果只有一个前提~$S_1$，那么只需满足 $\pheight{\pi} \le \pheight{\pi_1}$ 即可。对于 $\RightR{\Weakening}$ 和 $\LeftR{\Weakening}$ 规则，我们实际上有 $\pheight{\pi} = \pheight{\pi_1}$，但这不是必需的。

\begin{prop}\ollabel{prop:weak-G3c-adm}
\Log{G1c} 的规则 $\RightR{\Weakening}$ 和 $\LeftR{\Weakening}$ 在~\Log{G3c} 中是保持高度可容许的。
\end{prop}

\begin{proof}
  我们需要证明，如果 $\Gamma \Sequent \Delta$ 在~\Log{G3c} 中有一个高度为~$\pheight{\pi} = n$ 的推导~$\pi$，那么 $\Gamma \Sequent \Delta, !A$ 也有一个 \Log{G3c}-推导 $\pi'$，且 $\pheight{\pi'} = n$。我们对~$n$ 进行归纳证明。如果 $n=0$，那么 $\pi$ 就是一个公理 $\Gamma \Sequent \Delta$（即某个原子 $!C$ 同时出现在 $\Gamma$ 和~$\Delta$ 中），因而 $\Gamma \Sequent \Delta, !A$ 也是一条公理。如果 $n = \pheight{\pi} >0$，推导~$\pi$ 有最后一个推理。我们根据所使用的规则分情况讨论。我们只举一个例子：假设最后一个推理是 $\RightR{\land}$。那么 $\Delta = \Delta', !B \land !C$，并且该最后推理的两个直接子推导 $\pi_1$ 和~$\pi_2$ 的结束式分别是 $\Gamma \Sequent \Delta', !B$ 和 $\Gamma \Sequent \Delta', !C$。换言之，$\pi$ 看起来是这样的：
  \[
  \AxiomC{}
  \RightLabel{$\pi_1$}
  \Deduce$\Gamma \fCenter \Delta', !B$
  \AxiomC{}
  \RightLabel{$\pi_2$}
  \Deduce$\Gamma \fCenter \Delta', !C$
  \RightLabel{\RightR{\land}}
  \BinaryInf$\Gamma \fCenter \Delta', !B \land !C$
  \DisplayProof
  \]
  我们还知道 $n =
  \max(\pheight{\pi_1}, \pheight{\pi_2}) + 1$。因此 $\pheight{\pi_1} < n$ 且 $\pheight{\pi_2} < n$，归纳假设适用：存在推导 $\pi_1'$ 和~$\pi_2'$ 分别推出 $\Gamma \Sequent \Delta', !B, !A$ 和 $\Gamma \Sequent \Delta', !C, !A$。使用 $\RightR{\land}$ 规则，我们得到 $\Gamma
  \Sequent \Delta', !B \land !C, !A$ 的一个推导~$\pi'$：
  \[
  \AxiomC{}
  \RightLabel{$\pi_1'$}
  \Deduce$\Gamma \fCenter \Delta', !B, !A$
  \AxiomC{}
  \RightLabel{$\pi_2'$}
  \Deduce$\Gamma \fCenter \Delta', !C, !A$
  \RightLabel{\RightR{\land}}
  \BinaryInf$\Gamma \fCenter \Delta', !B \land !C, !A$
  \DisplayProof
  \]
  我们有 $\pheight{\pi'} = \max(\pheight{\pi_1'}, \pheight{\pi_2'}) + 1 = \max(\pheight{\pi_1}, \pheight{\pi_2}) + 1 = \pheight{\pi}$。
\end{proof}

\olref{prop:weak-G3c-adm} 极其有用。我们将为它的重复应用引入一个记号：如果 $\pi$ 是~$\Gamma \Sequent \Delta$ 的一个推导，那么 $\pi[\Pi \Sequent \Lambda]$ 表示对 $\pi$ 运用左弱化的可容许性（加入~$\Pi$ 中的所有公式）和右弱化的可容许性（加入~$\Lambda$ 中的所有公式）所得的结果。它是 $\Gamma, \Pi \Sequent \Delta, \Lambda$ 的一个推导。

\end{document}